\section{Introduction}\label{sec:intro}

Hierarchical task network (HTN) planning describes what an agent should do in terms of
\emph{tasks} rather than goals \parencite{erol:et:al:amai-96,ghallab:et:al:04,bercher:et:al:ijcai-19}.
Abstract (compound) tasks are refined by \emph{methods} into networks of simpler tasks until
only primitive tasks, i.e.\ actions, remain, and a plan is an executable arrangement of the
actions obtained in this way. When actions are deterministic, the computational complexity of
deciding whether a plan exists is well understood: it ranges from \NP{} for networks without
compound tasks to undecidability for unrestricted hierarchies, with a rich landscape in between
determined by the shape of the methods and by whether task networks are totally or only
partially ordered \parencite{erol:et:al:amai-96,geier:bercher:ijcai-11,alford:et:al:icaps-15}.

In many applications actions do not have a single predictable effect. In \emph{fully observable
nondeterministic} (FOND) planning an action has a set of possible outcomes; which one occurs is
determined by the environment, and the agent observes it before acting again. A plan must now
cope with all outcomes, and the question of what a plan \emph{is} becomes interesting. In
classical FOND planning, \textcite{cimatti:et:al:aij-03} distinguish \emph{weak} plans (some
execution reaches the goal), \emph{strong} plans (all executions reach the goal) and
\emph{strong cyclic} plans (executions may loop, but from every point a goal remains
reachable, so that under fair nondeterminism the goal is eventually reached).

For HTN planning there is a second dimension: \emph{when} the refinement decisions are taken.
\CBa{} studied solutions with \emph{fixed methods}: the planner first chooses a complete
refinement of the initial task network and only then executes the resulting primitive network,
either along a fixed order (\emph{linearisation-dependent}, LD) or with a policy that chooses
the next action depending on the outcomes so far (\emph{outcome-dependent}, OD). \CBb{}
introduced \emph{method-based policies}, which interleave refinement and execution, so that the
choice of a method may depend on everything observed before. Both papers determined the
complexity of plan existence for the standard structural classes of HTN problems, and
\DB{} filled one of the remaining entries.

\paragraph{Contributions.}
This article completes the picture. Table~\ref{tab:results} summarises the state of the art
after our work; the entries in bold are new. Concretely:

\begin{itemize}
  \item For fixed methods with strong outcome-dependent solutions we show that acyclic
    partially ordered problems are \EXPSPACE-complete, closing a gap between the known
    \NEXPTIME{} lower bound and the \EXPSPACE{} upper bound. We show that tail-recursive problems
    are undecidable, and so are regular problems in the sense of \textcite{bercher:et:al:ijcai-19},
    whose initial task network may contain several compound tasks. For regular problems whose
    initial network is regular, the \EXPSPACE-completeness established by \CBa{} and \DB{} is
    confirmed with a new, direct lower-bound proof.
  \item For method-based policies we show that totally ordered problems with an arbitrary
    hierarchy are \TWOEXP-complete, both for strong and for strong cyclic solutions; this row
    was missing from the table of \CBb{}. We also settle the strong cyclic entry for
    tail-recursive partially ordered problems (\TWOEXP-complete).
  \item We point out that the definition of strong cyclic solutions of \CBb{}, read literally,
    accepts policies that refine a recursive task forever, and that under this reading the
    lower bounds claimed for strong cyclic solutions do not hold. With the usual fairness
    condition of \textcite{cimatti:et:al:aij-03} they do, and all our strong cyclic results use
    that reading.
  \item Under this reading, strong cyclic plan existence for unrestricted problems is not
    semi-decidable, contrary to what one might expect from the other solution concepts. It is
    neither recursively enumerable nor co-recursively enumerable; in fact it is
    $\Sigma^1_1$-complete and therefore outside the arithmetical hierarchy.
  \item Several published proofs do not go through as written. The most consequential ones are
    the \PSPACE-hardness proof for primitive partially ordered problems with strong
    outcome-dependent solutions \parencite[Thm.~5.1]{chen:bercher:icaps-21}, whose reduction
    from $k$-in-a-row games can be won trivially, and the \EXPTIME-hardness proof for acyclic
    totally ordered problems with method-based policies
    \parencite[Thm.~1]{chen:bercher:icaps-22}, which lacks an acceptance check. Further proofs
    are outlines whose gadgets are not specified. We replace all of them by complete
    reductions.
\end{itemize}

\paragraph{Machine-checked proofs.}
Complexity results for planning formalisms rest on long reductions with many details, and the
experience recorded above shows that errors are easy to make and hard to spot. We have therefore
formalised the whole development in the proof assistant Lean~4
\parencite{moura:ullrich:cade-21} on top of Mathlib \parencite{mathlib:cpp-20}: the planning
formalism, all solution concepts, every membership and every hardness result of
Table~\ref{tab:results}, and the complexity classes themselves. The only result we use without
proof in Lean is Harel's theorem that recurrence of nondeterministic machines is
$\Sigma^1_1$-complete \parencite{harel:jacm-86}, together with a routine consequence for
two-stack machines; it enters only the $\Sigma^1_1$-hardness of one entry. Section~\ref{sec:formal}
describes the formalisation. The present article is written for readers who are interested in
the results and their proofs, not in the formalisation: we present the proofs as one would
write them on paper and keep references to the formal development to a minimum.

\paragraph{Outline.}
Section~\ref{sec:prelim} introduces the formalism and the structural classes.
Section~\ref{sec:solutions} defines the solution concepts, discusses the definition of strong
cyclic solutions, and relates the concepts to each other and to games.
Section~\ref{sec:fixed} treats fixed methods and Section~\ref{sec:policies} method-based
policies. Section~\ref{sec:unrestricted} is devoted to unrestricted partially ordered problems,
where the fair strong cyclic case turns out to be $\Sigma^1_1$-complete.
Section~\ref{sec:formal} describes the formalisation and Section~\ref{sec:conclusion} concludes.

\begin{table}[t]
\caption{The complexity of plan existence for FOND HTN problems. Rows are structural classes,
  each totally (\TO) or partially (\PO) ordered; columns are solution concepts: fixed methods
  (weak, strong linearisation-dependent, strong outcome-dependent) and method-based policies
  (weak, strong, strong cyclic in the fair reading). ``-c'' means complete. Bold entries are new
  in this article; entries marked $^\dagger$ were claimed before, but the published proof is
  incomplete or wrong and is replaced here. For regular problems, \REGtn{} requires the initial
  task network to be regular as well \parencite{chen:bercher:icaps-21}, while \REG{} does not
  \parencite{bercher:et:al:ijcai-19}. All entries are formally verified, except that the
  $\Sigma^1_1$-hardness in the last row relies on Harel's theorem.}
\label{tab:results}
\footnotesize
\setlength{\tabcolsep}{3pt}
\begin{tabularx}{\linewidth}{@{}ll YYY YYY@{}}
  \toprule
  & & \multicolumn{3}{c}{fixed methods} & \multicolumn{3}{c}{method-based policies} \\
  \cmidrule(lr){3-5}\cmidrule(l){6-8}
  class & & weak & strong LD & strong OD & weak & strong & strong cyclic \\
  \midrule
  \PRIM & \TO & \NP-c & in \PTIME & in \PTIME & \NP-c & in \PTIME & in \PTIME \\
        & \PO & \NP-c & \NP-c & \PSPACE-c$^\dagger$ & \NP-c & \PSPACE-c$^\dagger$ & \PSPACE-c$^\dagger$ \\
  \ACYC & \TO & \PSPACE-c & \PSPACE-c & \PSPACE-c & \PSPACE-c & \EXPTIME-c$^\dagger$ & \EXPTIME-c$^\dagger$ \\
        & \PO & \NEXPTIME-c & \NEXPTIME-c & \textbf{\EXPSPACE-c} & \NEXPTIME-c & \EXPSPACE-c$^\dagger$ & \EXPSPACE-c$^\dagger$ \\
  \REG, \REGtn & \TO & \PSPACE-c & \PSPACE-c & \PSPACE-c & \PSPACE-c & \EXPTIME-c$^\dagger$ & \EXPTIME-c$^\dagger$ \\
  \REGtn & \PO & \PSPACE-c & \PSPACE-c & \EXPSPACE-c$^\dagger$ & \PSPACE-c & \EXPTIME-c$^\dagger$ & \EXPTIME-c$^\dagger$ \\
  \REG  & \PO & \PSPACE-c & \PSPACE-c & \textbf{undecidable} & \PSPACE-c & \EXPTIME-c$^\dagger$ & \EXPTIME-c$^\dagger$ \\
  \TR   & \TO & \PSPACE-c & \PSPACE-c & \PSPACE-c & \PSPACE-c & \EXPTIME-c$^\dagger$ & \EXPTIME-c$^\dagger$ \\
        & \PO & \EXPSPACE-c & \EXPSPACE-c & \textbf{undecidable} & \EXPSPACE-c & \TWOEXP-c$^\dagger$ & \textbf{\TWOEXP-c} \\
  \ARB  & \TO & \EXPTIME-c & \EXPTIME-c & \EXPTIME-c & \EXPTIME-c & \textbf{\TWOEXP-c} & \textbf{\TWOEXP-c} \\
        & \PO & undec., \RE & undec., \RE & undec., \RE & undec., \RE & undec., \RE & \textbf{$\Sigma^1_1$-c} \\
  \bottomrule
\end{tabularx}
\end{table}
